\chapter{Обробка даних.}

Обробка даних (Data Processing) — це сукупність методів і процедур, спрямованих на підготовку 
даних до подальшого аналізу та побудови моделей машинного навчання. Основними етапами обробки 
є очищення даних, усунення пропущених значень і дублікатів, виявлення викидів, перетворення 
типів даних, а також масштабування та нормалізація ознак. Якість виконання цих етапів 
безпосередньо впливає на точність, швидкість і надійність моделей машинного навчання, 
оскільки алгоритми навчаються саме на підготовлених даних. «Сирі» дані часто містять 
пропущені значення, дублікати, помилки введення, викиди, некоректні формати та 
неоднорідні масштаби ознак, що може призводити до погіршення якості аналізу та зниження 
точності прогнозування.


\section{Pandas}

Існує ряд бібліотек Python, які обробляють реляційні дані, зазвичай написані як 
інтерфейси до декількох різних систем управління реляційними базами даних (таких як 
PostgreSQL, MySQL та інші). 

Примітка: слід зазначити, що програмне забезпечення, таке 
як PostreSQL і MySQL, правильніше називати системами управління реляційними базами даних 
(RDBMS), а не базами даних. База даних — це фактичні таблиці та записи, що визначають 
фактичну сукупність даних.

Однак у цьому курсі ми будемо працювати з реляційними даними переважно за допомогою 
двох бібліотек: Pandas і SQLite. Це особливо прості бібліотеки, якщо говорити про 
реальні бази даних: Pandas, безумовно, не є реальною системою реляційних баз даних 
(хоча вона надає функції, що віддзеркалюють деякі їхні можливості), тоді як SQLite 
є «реальною» RDBMS, але надзвичайно простою, без стандартної архітектури 
клієнт/сервер, яка є практично у всіх реальних виробничих базах даних. Проте для 
багатьох задач у галузі науки про дані вони будуть достатніми, тому ми зосередимося на них тут.

Ми вже коротко розглянули Panda, коли обговорювали збір даних, і вона виявилася однією 
з найкорисніших бібліотек Python для науки про дані. Як ми вже згадували вище (але 
будемо повторювати цей факт багато разів), Pandas — це не бібліотека реляційних баз 
даних, а бібліотека «даних-рам» (DataFrame). Ви можете уявити собі фрейм даних як 2D-масив, за 
винятком того, що записи у фреймі даних можуть бути будь-яким типом об'єкта Python 
(і мати змішані типи в масиві), а рядки/стовпці можуть мати «мітки» замість 
цілочисельних індексів, як у стандартному масиві.

Давайте подивимося, як спочатку створити фрейм даних у Pandas, який відображає таблицю Person 
із записами про деяких людей.

\begin{lstlisting}
import pandas as pd

df = pd.DataFrame([(1, 'Ivanov', 'Ivan'), 
                   (2, 'Petrov', 'Petro'),
                   (3, 'Kosovan', 'Vasyl'), 
                   (4, 'Cherevko', 'Igor'),
                   (5, 'Kravtsov', 'Kyryl'),
                   (6, 'Rarenko', 'Anastasiya'),], 
                  columns=["id", "last_name", "first_name"])
df
\end{lstlisting}

Результат:

\begin{console}
    id	last_name	first_name
0	1	Ivanov	Ivan
1	2	Petrov	Petro
2	3	Kosovan	Vasyl
3	4	Cherevko	Igor
4	5	Kravtsov	Kyryl
5	6	Rarenko	Anastasiya
\end{console}

Ви, напевно, пам'ятаєте, що таке первинний ключ в базі даних. Якщо ні, можна коротко згадати, що 
це одне (або декілька) полів (або стовпців) у таблиці бази даних, які унікально ідентифікують 
кожен її рядок.

«Індекс» для Pandas насправді означає щось більше, ніж «первинний ключ» у таблиці бази 
даних (хоча б з тим винятком, що можливі дублікати записів). 

Тобто індекс (якщо він зроблений правильно, без дублікатів) є ідентифікатором для кожного 
рядка в базі даних. Ми можемо встановити індекс для одного з існуючих стовпців за допомогою 
виклику \verb|.set_index()|.

Але тут потрібно бути дуже обережним з одним моментом. За замовчуванням більшість операцій 
Pandas, таких як \verb|.set_index()| та багато інших, не виконуються на місці. Тобто, хоча 
виклик \verb|df.set_index(«id»)| вище повертає копію фрейму даних df з індексом, встановленим 
для стовпця id (пам'ятайте, що Jupyter notebook відображає значення, яке повертається 
в останньому рядку комірки), оригінальний об'єкт \verb|df| насправді тут не змінюється.

Якщо ми хочемо фактично змінити сам об'єкт \verb|df|, потрібно використовувати прапор 
\verb|inplace=True| для цих функцій (або присвоїти оригінальний об'єкт результату функції, 
але це не так чітко).

\begin{lstlisting}
df.set_index("id", inplace=True)
df
\end{lstlisting}

Результат:

\begin{console}
	last_name	first_name
id		
1	Ivanov	Ivan
2	Petrov	Petro
3	Kosovan	Vasyl
4	Cherevko	Igor
5	Kravtsov	Kyryl
6	Rarenko	Anastasiya
\end{console}

Series є одновимірною структурою, тоді як DataFrame є двовимірною. Наприклад, один стовпець таблиці: \verb|df['last_name']| повертає об'єкт Series.

Якщо ви захочете завантажити дані із CSV-файлу у DataFrame, це можна зробити наступним чином:

\begin{lstlisting}
import pandas as pd

data = pd.read_csv('users.csv')
\end{lstlisting}

Для отримання загальної інформації про структуру DataFrame використовується метод info(). Він показує кількість рядків, назви стовпців, кількість непорожніх значень і типи даних. Наприклад, за його допомогою можна швидко визначити, у яких стовпцях є пропущені значення.

\begin{lstlisting}
df.info()
\end{lstlisting}

Результат:

\begin{console}
<class 'pandas.DataFrame'>
RangeIndex: 6 entries, 1 to 6
Data columns (total 2 columns):
 #   Column      Non-Null Count  Dtype
---  ------      --------------  -----
 0   last_name   6 non-null      str  
 1   first_name  6 non-null      str  
dtypes: str(2)
memory usage: 303.0 bytes
\end{console}

Для отримання розміру таблиці використовують атрибут shape: \verb|df.shape|. Він повертає кортеж із кількістю рядків і стовпців.

Для статистичного опису числових даних використовується метод \verb|df.describe()|. Він автоматично обчислює основні статистичні характеристики числових стовпців, зокрема кількість значень, середнє значення, стандартне відхилення, мінімальне та максимальне значення, а також квартилі.

\textbf{Розглянемо іще деякі методи для роботи із фреймами в Pandas.}

Для перегляду перших записів використовується метод \verb|head()|. Наприклад, перші три рядки:

\begin{lstlisting}
df.head(3)
\end{lstlisting}

Результат:

\begin{console}
	last_name	first_name
id		
1	Ivanov	Ivan
2	Petrov	Petro
3	Kosovan	Vasyl
\end{console}

Для перегляду останніх записів використовується метод \verb|tail()|. Наприклад, останні три рядки:

\begin{lstlisting}
df.tail(3)
\end{lstlisting}

Результат:

\begin{console}
   last_name	first_name
id		
4	Cherevko	Igor
5	Kravtsov	Kyryl
6	Rarenko	Anastasiya
\end{console}

Доступ до окремих елементів можна отримати за допомогою нотації \verb|.loc[row, column]|, 
де row позначає індекс, який ви шукаєте, а column - назву стовпця.

\begin{lstlisting}
df.loc[1, "last_name"]
\end{lstlisting}

Результат:

\begin{console}
'Ivanov'
\end{console}

Якщо ми хочемо отримати доступ до всіх прізвищ (або до всіх елементів у певному рядку), 
ми використовуємо символ-замінник \verb|:|.

\begin{lstlisting}
df.loc[:, "last_name"]
\end{lstlisting}

Результат:

\begin{console}
id
1      Ivanov
2      Petrov
3     Kosovan
4    Cherevko
5    Kravtsov
6     Rarenko
Name: last_name, dtype: str
\end{console}

Ми можемо передати список потрібних стовпців, щоб отримати об’єкт типу DataFrame.

\begin{lstlisting}
df.loc[:, ["last_name"]]
\end{lstlisting}

Результат:

\begin{console}
    last_name
id	
1	Ivanov
2	Petrov
3	Kosovan
4	Cherevko
5	Kravtsov
6	Rarenko
\end{console}

Те саме можна зробити з індексами рядків.

\begin{lstlisting}
df.loc[[1,2],:]
\end{lstlisting}

Результат:

\begin{console}
	last_name	first_name
id		
1	Ivanov	Ivan
2	Petrov	Petro
\end{console}



Крім того, ми можемо використовувати `.loc` для зміни вмісту існуючих записів.

\begin{lstlisting}
df.loc[6,"last_name"] = "Vikovan"
df
\end{lstlisting}

Результат:

\begin{console}
	last_name	first_name
id		
1	Ivanov	Ivan
2	Petrov	Petro
3	Kosovan	Vasyl
4	Cherevko	Igor
5	Kravtsov	Kyryl
6	Vikovan	Anastasiya
\end{console}




Ми можемо навіть додати додаткові рядки/стовпці.

\begin{lstlisting}
df.loc[7,:] = ('Vikovan', 'Valentyn')
df
\end{lstlisting}

Результат:

\begin{console}
last_name	first_name
id		
1	Ivanov	Ivan
2	Petrov	Petro
3	Kosovan	Vasyl
4	Cherevko	Igor
5	Kravtsov	Kyryl
6	Vikovan	Anastasiya
7	Vikovan	Valentyn
\end{console}


\ding{43} Наостанок пам’ятайте, що `.loc` завжди індексує на основі «індексу» (тобто, 
фактично, первинного ключа) фрейму даних 
разом з назвою стовпця. Якщо ви хочете отримати доступ на основі позиційного 
індексу (тобто, використовуючи лічильники з нульовим індексом 
як для рядків, так і для стовпців), ви можете скористатися властивістю \verb|.iloc|.

\begin{lstlisting}
df.iloc[4,1]
\end{lstlisting}

Результат:

\begin{console}
'Kyryl'
\end{console}

Рядки DataFrame можна відфільтрувтаи за певною умовою:

\begin{lstlisting}
filtered = df[df['first_name'].str.contains('I')]
display(filtered)
\end{lstlisting}

Результат:

\begin{console}
	last_name	first_name
id		
1	Ivanov	Ivan
4	Cherevko	Igor
\end{console}

Можна також комбінувати декілька умов.

Новий стовпець можна створити за допомогою оператора присвоєння:

\begin{lstlisting}
df['user_name'] = df['last_name'] + df['first_name']
display(df)
\end{lstlisting}

Результат:

\begin{console}
	last_name	first_name	user_name
id			
1	Ivanov	Ivan	IvanovIvan
2	Petrov	Petro	PetrovPetro
3	Kosovan	Vasyl	KosovanVasyl
4	Cherevko	Igor	CherevkoIgor
5	Kravtsov	Kyryl	KravtsovKyryl
6	Vikovan	Anastasiya	VikovanAnastasiya
7	Vikovan	Valentyn	VikovanValentyn
\end{console}

Для видалення стовпця використовується метод drop():

\begin{lstlisting}
df = df.drop('user_name', axis=1)
\end{lstlisting}

Параметр axis=1 означає, що видаляється стовпець. Для видалення рядка можна використати його індекс та параметр axis=0.

Для перейменування використовують метод rename():

\begin{lstlisting}
df = df.rename(columns={
    'last_name': 'Last_Name',
})
\end{lstlisting}

Для сортування використовується метод \verb|sort_values()|:

\begin{lstlisting}
df = df.sort_values('first_name')
\end{lstlisting}

\section{Методи очищення даних та обробка відсутніх значень.}

У реальному світі необроблені дані рідко бувають ідеальними. Вони часто містять помилки, 
невідповідності та відсутні значення. 

Навіть найкращий алгоритм чи штучний інтелект не зможе побудувати хорошу модель, якщо навчальні дані 
містять помилки або пропуски. Наприклад, прогноз погоди буде ускладненим, якщо серед показників 
термометра за попередні дні затесалися 
такі дані як -500 градусів за Цельсієм, чи просто багато пропусків.

Очищення даних — це процес виявлення та виправлення 
цих проблем для поліпшення якості даних.

\begin{itemize}
\item \textit{Дублікати}: це повторювані записи в наборі даних. Дублікати рядків можуть спотворити 
ваш аналіз або модель, особливо якщо вони представляють одну й ту саму подію або об'єкт.
\item \textit{Неправильні типи даних}: дані, які введені неправильно (наприклад, числа, збережені як 
рядки), можуть призвести до помилок в аналізі та навчанні моделі.
\item \textit{Випадкові значення}: екстремальні значення, які значно відрізняються від решти даних. 
Вони можуть спотворити результати статистичного аналізу або моделей машинного навчання.
\end{itemize}

Відсутність даних може відбуватися з різних причин, таких як помилки введення даних, 
несправності обладнання або проблеми з конфіденційністю. Важливо обережно поводитися 
з відсутніми значеннями, оскільки неправильне поводження може призвести до упереджених 
результатів.

Типи відсутніх даних:

\begin{itemize}
\item \textit{MCAR} (Missing Completely at Random, повністю випадкова відсутність): ймовірність 
відсутності даних не залежить від самих даних.
\item \textit{MAR} (відсутність випадкова): ймовірність відсутності даних пов'язана з іншими 
спостережуваними даними, але не з самими відсутніми даними.
\item \textit{MNAR} (відсутність невипадкова): відсутність пов'язана з самими неспостережуваними 
даними.
\end{itemize}



\subsection{Видалення відсутніх даних.}

Видалення рядків або стовпців з відсутніми значеннями є найпростішим підходом, але може 
призвести до втрати цінної інформації.

У наступному прикладі метод \verb|df.dropna()| видаляє всі рядки, у яких є хоча б одне відсутнє 
значення, а метод \verb|df.dropna(subset=['feature1', 'feature2'])| видаляє 
рядки, у яких є хоча б одне відсутнє значення в числових стовпцях feature1 та feature2.

\begin{lstlisting}
import pandas as pd

df = pd.read_csv('../resources/data.csv')
print(df, '\n')

df_cleaned = df.dropna()  
df_feature_cleaned = df.dropna(subset=['feature1', 'feature2'])
print(df_cleaned, '\n')
print(df_feature_cleaned, '\n')
\end{lstlisting}

Результат:

\begin{console}
   feature1  feature2 categorical_column  target
0       1.2       3.2                  A      10
1       2.3       NaN                  B      15
2       3.1       1.1                  C       7
3       4.7       4.8                  A      18
4       NaN       2.7                  B      11
5       5.1       4.2                  C      19
6       2.8       NaN                  A      13
7       4.1       3.9                  B      16
8       5.3       4.1                NaN      14
9       3.3       2.1                  C      12 

   feature1  feature2 categorical_column  target
0       1.2       3.2                  A      10
2       3.1       1.1                  C       7
3       4.7       4.8                  A      18
5       5.1       4.2                  C      19
7       4.1       3.9                  B      16
9       3.3       2.1                  C      12 

   feature1  feature2 categorical_column  target
0       1.2       3.2                  A      10
2       3.1       1.1                  C       7
3       4.7       4.8                  A      18
5       5.1       4.2                  C      19
7       4.1       3.9                  B      16
8       5.3       4.1                NaN      14
9       3.3       2.1                  C      12 
\end{console}

За допомогою методу \verb|df.isna()| можна виявити пропущенні значення у DataFrame. Щоб отримати кількість пропусків у кожному стовпці, зручно використати \verb|df.isna().sum()|

\ding{43} Видалення рядків доцільне тоді, коли кількість пропущених записів невелика і їхнє видалення не призведе до істотної втрати інформації. Наприклад, якщо з 100 000 записів лише 20 містять пропущені значення, їх можна видалити без значного впливу на розмір вибірки. Водночас якщо пропуски присутні у значній частині даних або стосуються важливої ознаки, краще розглянути імпутацію. Особливо обережно потрібно працювати з пропусками у невеликих наборах даних, де видалення навіть декількох відсотків записів може суттєво змінити результати аналізу.


\subsection{Заповнення середнім.} 

Імпутація даних (Data Imputation) — це процес заміни пропущених значень оціненими значеннями.

При заповненні середнім (Mean Imputation) ми замінюємо відсутні значення середнім значенням стовпця.

\begin{lstlisting}
df['feature1'].fillna(df['feature1'].mean(), inplace=True)
df['feature2'].fillna(df['feature2'].mean(), inplace=True)
df
\end{lstlisting}

Результат:

\begin{console}
feature1	feature2	categorical_column	target
0	1.200000	3.2000	A	10
1	2.300000	3.2625	B	15
2	3.100000	1.1000	C	7
3	4.700000	4.8000	A	18
4	3.544444	2.7000	B	11
5	5.100000	4.2000	C	19
6	2.800000	3.2625	A	13
7	4.100000	3.9000	B	16
8	5.300000	4.1000	NaN	14
9	3.300000	2.1000	C	12
\end{console}

\subsection{Заміщення медіанним значенням.} 

При заміщенні медіанним значенням (Median Imputation) ми заміщуємо відсутні значення медіанним значенням. 
Цей метод є корисним, коли дані є викривленими.

\begin{lstlisting}
df = pd.read_csv('../resources/data.csv')
df['feature1'].fillna(df['feature1'].median(), inplace=True)
df
\end{lstlisting}

Результат:

\begin{console}
	feature1	feature2	categorical_column	target
0	1.2	3.2	A	10
1	2.3	NaN	B	15
2	3.1	1.1	C	7
3	4.7	4.8	A	18
4	3.3	2.7	B	11
5	5.1	4.2	C	19
6	2.8	NaN	A	13
7	4.1	3.9	B	16
8	5.3	4.1	NaN	14
9	3.3	2.1	C	12
\end{console}

\subsection{Заміщення модою.} 

Заміщення модою (Mode Imputation) часто використовується для категоріальних даних.

\begin{lstlisting}
df['categorical_column'].fillna(df['categorical_column'].mode()[0], inplace=True)
df
\end{lstlisting}

Результат:

\begin{console}
feature1	feature2	categorical_column	target
0	1.2	3.2	A	10
1	2.3	NaN	B	15
2	3.1	1.1	C	7
3	4.7	4.8	A	18
4	3.3	2.7	B	11
5	5.1	4.2	C	19
6	2.8	NaN	A	13
7	4.1	3.9	B	16
8	5.3	4.1	A	14
9	3.3	2.1	C	12
\end{console}

\ding{43} Інші поширені методи заповнення даних:

\textit{Імпутування K-найближчих сусідів (KNN):} оцінює відсутні значення на основі значень найближчих 
сусідів.

\textit{Багатовимірне імпутування (підстановка):} враховує взаємозв'язки між ознаками для імпутування 
відсутніх значень.

\textit{Створення індикаторних змінних:} інший підхід полягає у створенні бінарної індикаторної 
змінної, яка позначає наявність відсутніх даних.


\section{Дублікати.}

Дублікати — це записи, які повністю або частково повторюють інші записи у наборі даних. Дублікати можуть виникати під час повторного завантаження даних, об'єднання декількох джерел або помилкового введення інформації.

Для перевірки дубльованих рядків використовується метод duplicated():

\begin{lstlisting}
df.duplicated()
\end{lstlisting}

Він повертає True для рядків, які є дубльованими.

Щоб отримати самі дублікати:

\begin{lstlisting}
df[df.duplicated()]
\end{lstlisting}

Кількість дублікатів можна визначити так:

\begin{lstlisting}
df.duplicated().sum()
\end{lstlisting}

Для видалення дубльованих рядків використовується метод \verb|drop_duplicates()|:

\begin{lstlisting}
df = df.drop_duplicates()
\end{lstlisting}

Можна визначати дублікати лише за певними стовпцями.

\begin{lstlisting}
df = df.drop_duplicates(subset=['feature1', 'feature2'])
\end{lstlisting}



\section{Масштабування та нормалізація даних.}

Масштабування та нормалізація даних (Data Scaling and Normalization) є важливими етапами попередньої обробки, які можуть суттєво 
вплинути на ефективність ваших моделей.

Масштабування даних гарантує, що всі ознаки однаково впливають на процес навчання моделі. 
Коли ознаки мають дуже різні діапазони, ті, що мають більші діапазони, можуть домінувати в 
обчисленнях відстані або кроках градієнта, що призводить до неоптимальних моделей.

Наприклад, якщо Age має значення від 18 до 80, а Salary — від 20 000 до 200 000, то без масштабування більші за величиною значення Salary можуть непропорційно впливати на роботу моделі.

\textit{Алгоритми на основі відстані:} В алгоритмах, таких як K-найближчі сусіди (KNN) або машини 
опорних векторів (SVM), відстань між точками даних має вирішальне значення. Якщо одна ознака 
має набагато більший діапазон, ніж інші, вона непропорційно впливатиме на метрику відстані, 
що може спотворити результати.

\textit{Градієнтний спуск:} Алгоритми, такі як лінійна регресія та нейронні мережі, використовують 
градієнтний спуск для мінімізації функції вартості. Якщо ознаки не масштабуються, алгоритм 
градієнтного спуску може потребувати більше часу для збіжності або може збігатися до 
неоптимального рішення через нерівномірні кроки, зроблені в просторі параметрів.

Масштабування - це загальна назва перетворень, які змінюють масштаб числових ознак. Воно може включати як приведення значень до певного інтервалу, так і стандартизацію.

Нормалізація у контексті підготовки даних часто означає перетворення значень до визначеного діапазону, наприклад від 0 до 1.

Тому мінімально-максимальна нормалізація є одним із видів масштабування, тоді як стандартизація є іншим способом масштабування.

\subsection{Викиди.}

Викид (Outlier) це спостереження, значення якого істотно відрізняється від більшості інших значень у вибірці.

Наприклад, якщо зарплати працівників знаходяться переважно в межах від 20 000 до 80 000 грн, значення 1 000 000 грн може бути потенційним викидом.

Викид не обов'язково є помилкою. Він може бути результатом помилки вимірювання або введення даних, але також може відображати реальне рідкісне явище. Тому викиди не слід автоматично видаляти.

Одним із найпоширеніших статистичних методів є правило міжквартильного розмаху (IQR). Значення вважається потенційним викидом, якщо воно знаходиться нижче $Q_{1} - 1.5 \cdot IQR$ або вище $Q_{3} + 1.5 \cdot IQR$, де $Q_{1}$ і $Q_{3}$ - перший та третій квартилі, а $IQR = Q_{3} - Q_{1}$. 

Також викиди можна виявляти графічно за допомогою boxplot, гістограм або scatter plot, а для складніших багатовимірних даних — за допомогою алгоритмів кластеризації та спеціальних методів виявлення аномалій. про це ми поговоримо в наступних розділах.

Викиди можуть суттєво змінювати статистичні характеристики вибірки. Наприклад, один надзвичайно великий показник може значно збільшити середнє значення.

Вони також можуть впливати на параметри деяких моделей машинного навчання. Особливо чутливими до викидів можуть бути методи, які використовують відстані або мінімізують квадрати помилок, наприклад лінійна регресія та k-середніх.

Однак видалення викидів також може бути небезпечним. Якщо викид є реальним явищем, його видалення означатиме втрату важливої інформації. Тому спочатку потрібно з'ясувати природу аномального значення.

\subsection{Мінімально-максимальне масштабування.}

Мінімально-максимальне масштабування (min-max normalization) перемасштабовує дані до фіксованого діапазону, зазвичай \verb|[0, 1]|. Ця техніка особливо корисна, 
коли ви знаєте, що ваші дані відповідають розподілу з чіткими верхньою та нижньою межами.

$$X' = \frac{X-X_{min}}{X_{max}-X_{min}}$$

\begin{itemize}
\item Використовуйте мінімально-максимальне масштабування, коли ви хочете, щоб усі ваші ознаки 
мали однаковий масштаб (наприклад, в алгоритмах, таких як KNN, SVM).
\item Особливо корисно, коли дані розподілені по відомому і фіксованому діапазону.
\end{itemize}

\begin{lstlisting}
from sklearn.preprocessing import MinMaxScaler
import pandas as pd

data = {'feature': [10, 50, 100, 150, 200]}
df = pd.DataFrame(data)

scaler = MinMaxScaler()
scaled_data = scaler.fit_transform(df)

print(scaled_data)
\end{lstlisting}

Результат:

\begin{console}
[[0.        ]
 [0.21052632]
 [0.47368421]
 [0.73684211]
 [1.        ]]
\end{console}

Ми можемо вибрати окремі стовпці для масштабування через 
\verb|numerical_scaled = scaler.fit_transform(numerical_cols)|. 

В наступному прикладі ми спочатку 
вибираємо числові стовпці для масштабування через 
\verb|numerical_cols = df[['feature1', 'feature2']]|, потім розділяємо стовпці категорій та цільові 
стовпці
\verb|categorical_cols = df[['categorical_column', 'target']]|,
інціалізуємо масштабування 
\verb|scaler = MinMaxScaler()|,
підготовуємо та перетворюємо числові стовпці
\verb|numerical_scaled = scaler.fit_transform(numerical_cols)|,
перетворюємо нормалізовані дані назад у DataFrame 
\verb|numerical_scaled_df = pd.DataFrame(numerical_scaled, columns=numerical_cols.columns)|,
та об’єднуємо нормалізовані числові дані з категоріальними стовпцями
\verb|df_scaled = pd.concat([numerical_scaled_df, categorical_cols.reset_index(drop=True)], axis=1)|.

\begin{lstlisting}
df = pd.read_csv('../resources/data.csv')

numerical_cols = df[['feature1', 'feature2']]

categorical_cols = df[['categorical_column', 'target']]

from sklearn.preprocessing import MinMaxScaler

scaler = MinMaxScaler()

numerical_scaled = scaler.fit_transform(numerical_cols)

numerical_scaled_df = pd.DataFrame(numerical_scaled, columns=numerical_cols.columns)

df_scaled = pd.concat([numerical_scaled_df, categorical_cols.reset_index(drop=True)], axis=1)

print(df_scaled)
\end{lstlisting}

Результат:

\begin{console}
   feature1  feature2 categorical_column  target
0  0.000000  0.567568                  A      10
1  0.268293       NaN                  B      15
2  0.463415  0.000000                  C       7
3  0.853659  1.000000                  A      18
4       NaN  0.432432                  B      11
5  0.951220  0.837838                  C      19
6  0.390244       NaN                  A      13
7  0.707317  0.756757                  B      16
8  1.000000  0.810811                NaN      14
9  0.512195  0.270270                  C      12
\end{console}

\ding{43} Основним недоліком Min-Max нормалізації є її чутливість до викидів. Якщо одна ознака містить надзвичайно велике або мале значення, $x_{min}$ або $x_{max}$ можуть сильно змінитися, внаслідок чого більшість нормалізованих значень опиниться у вузькому діапазоні. Крім того, якщо у майбутніх даних з'являться значення, що виходять за межі початкового діапазону, вони можуть отримати значення менше 0 або більше 1.


\subsection{Нормалізація (стандартизація) за Z-балом.}

Нормалізація (стандартизація) за Z-балом (Z-Score Normalization or standardization) перемасштабовує дані так, щоб середнє значення дорівнювало 0, а стандартне відхилення — 1. 
Ця техніка корисна, коли дані мають гаусівський (нормальний) розподіл, але не обов'язково обмежені.

$$X' = \frac{X-\nu}{\sigma}$$

\begin{itemize}
\item Використовуйте нормалізацію за Z-балом, коли ваші дані мають нормальний розподіл.
\item Це стандартний вибір для багатьох моделей машинного навчання, особливо тих, що передбачають нормальний розподіл даних.
\end{itemize}

\begin{lstlisting}
from sklearn.preprocessing import StandardScaler
import pandas as pd

data = {'feature': [10, 50, 100, 150, 200]}
df = pd.DataFrame(data)

scaler = StandardScaler()
standardized_data = scaler.fit_transform(df)

print(standardized_data)
\end{lstlisting}

Результат:

\begin{console}
[[-1.35411306]
 [-0.76536825]
 [-0.02943724]
 [ 0.70649377]
 [ 1.44242478]]
\end{console}

Нормалізація дозволяє привести ознаки до єдиного діапазону, що особливо корисно для алгоритмів, які використовують відстані.

Вона запобігає ситуації, коли ознака з великими числовими значеннями домінує над ознаками з меншим масштабом. Крім того, нормалізація може покращити стабільність та швидкість оптимізації деяких моделей.


\subsection{Надійне масштабування.}

Надійне масштабування (Robust Scaling) використовує медіану та міжквартильний розмах (IQR) для масштабування, що робить його 
надійним щодо винятків. На відміну від мінімально-максимального масштабування та 
нормалізації Z-балу, надійне масштабування менше піддається впливу екстремальних 
значень у даних.

$$X'=\frac{X - \text{медіана}(X)} {\text{IQR}(X)}$$

Де $\text{IQR}(X)$ — діапазон між 25-м процентилем ($Q1$) і 75-м процентилем ($Q3$).

Використовуйте надійне масштабування, якщо ваші дані містять винятки або мають негаусівський розподіл.
Воно ідеально підходить для ситуацій, коли дані є асиметричними або містять екстремальні значення, 
які не повинні сильно впливати на масштабування.

\begin{lstlisting}
from sklearn.preprocessing import RobustScaler
import pandas as pd

data = {'feature': [10, 50, 90, 100, 110, 150, 160, 200]}
df = pd.DataFrame(data)

scaler = RobustScaler()
robust_scaled_data = scaler.fit_transform(df)

print(robust_scaled_data)
\end{lstlisting}

Результат:

\begin{console}
[[-1.42857143]
 [-0.85714286]
 [-0.28571429]
 [-0.14285714]
 [ 0.        ]
 [ 0.57142857]
 [ 0.71428571]
 [ 1.28571429]
 [34.14285714]]
\end{console}

\ding{43} Для яких випадків підходить той чи інший метод масштабування?

\textit{Мінімально-максимальне масштабування}: використовується, коли відомі межі даних або коли 
алгоритм вимагає даних у певному діапазоні.

\textit{Нормалізація за Z-балом}: підходить для нормально розподілених даних або коли потрібно 
стандартизувати ознаки, щоб вони мали однакову важливість.

\textit{Надійне (Robust) масштабування}: найкраще підходить для роботи з даними, що м
істять винятки.



\section{Вступ до конструювання ознак.}

Інженерія ознак (Feature Engineering) є важливим етапом у процесі попередньої обробки даних. Вона передбачає перетворення 
необроблених даних на ознаки, які краще відображають основну проблему для моделей прогнозування. 
Якісна інженерія ознак може значно покращити ефективність моделей машинного навчання, надаючи їм 
найбільш релевантні та інформативні вхідні дані.

Для прикладу, маючи такі фізичні дані про людину як її вага та зріст, ми можемо 
сконструювати нову ознаку - "індекс маси тіла", що буде значно користнішим для нашої моделі,
ніж дві інші ознаки кожна окремо. 

\textit{Інженерія ознак} - це процес відбору, модифікації або створення нових ознак з необроблених даних 
для покращення ефективності моделей машинного навчання. Метою є створення ознак, які ефективніше 
відображають основні закономірності в даних, ніж вихідні необроблені дані. Правильна інженерія 
ознак може призвести до спрощення моделей, підвищення точності та скорочення часу навчання.

\subsection{Створення нових ознак.}

Створення нових ознак це виведення нових змінних з існуючих, які краще відображають основну структуру даних. 
Це може включати математичні перетворення, агрегації або операції, специфічні для певної області.

В наступному прикладі, маючи такі дані як вага та ріст людини, ми створюємо нову ознаку: 
BMI (body mass index, індекс маси тіла).

\begin{lstlisting}
import pandas as pd

data = {
    'height': [1.60, 1.75, 1.82, 1.90],
    'weight': [55, 80, 72, 90]
}
df = pd.DataFrame(data)

df

df['BMI'] = df['weight'] / (df['height'] ** 2)

df
\end{lstlisting}

Результат:

\begin{console}
	height	weight	BMI
0	1.60	55	21.484375
1	1.75	80	26.122449
2	1.82	72	21.736505
3	1.90	90	24.930748
\end{console}


\subsection{Кодування категоріальних змінних.}

Алгоритми машинного навчання зазвичай вимагають числових даних. Тому категоріальні змінні необхідно 
перетворити в числовий формат. Існує кілька методів кодування категоріальних змінних:

Кодування \textit{One-hot} (унітарне кодування) перетворює кожну категорію в новий бінарний стовпець. 
Цей метод підходить для 
випадків, коли кількість категорій обмежена.


\begin{lstlisting}
data = {'color': ['red', 'blue', 'green']}
df = pd.DataFrame(data)

# One-Hot Encoding
df_encoded = pd.get_dummies(df, columns=['color'])

print(df_encoded)
\end{lstlisting}

Результат:

\begin{console}
   color_blue  color_green  color_red
0           0            0          1
1           1            0          0
2           0            1          0
\end{console}

\textit{Кодування міток} (Label Encoding) присвоює кожної категорії унікальне ціле число. Це простий 
метод, 
але він передбачає порядкові відносини між категоріями, що може бути недоречним у деяких випадках.


\begin{lstlisting}
from sklearn.preprocessing import LabelEncoder

data = {'color': ['red', 'blue', 'green']}
df = pd.DataFrame(data)

# Label Encoding
encoder = LabelEncoder()
df['color_encoded'] = encoder.fit_transform(df['color'])

df
\end{lstlisting}

Результат:

\begin{console}
	color	color_encoded
0	red	2
1	blue	0
2	green	1
\end{console}

\textit{Цільове кодування} (Target Encoding) передбачає заміну кожної категорії середнім значенням 
цільової змінної для цієї категорії. Ця техніка є корисною, коли категоріальна змінна має 
велику кількість рівнів.

\begin{lstlisting}
data = {
    'color': ['red', 'blue', 'green', 'red', 'blue', 'green'],
    'target': [1, 0, 0, 1, 1, 0]
}
df = pd.DataFrame(data)

# Target Encoding
df['color_encoded'] = df.groupby('color')['target'].transform('mean')

print(df)
\end{lstlisting}

Результат:

\begin{console}
   color  target  color_encoded
0    red       1            1.0
1   blue       0            0.5
2  green       0            0.0
3    red       1            1.0
4   blue       1            0.5
5  green       0            0.0
\end{console}

\subsection{Функція взаємодії.}

Функція взаємодії створюються шляхом поєднання двох або більше змінних для відображення 
взаємодії між ними. Ці функції можуть виявити взаємозв'язки, які не є очевидними при 
розгляді змінних окремо.

Приклад:

Припустимо, що ви маєте такі дані, як вік і дохід користувачів. Термін взаємодії 
можна створити шляхом множення цих двох властивостей, щоб відобразити, як дохід 
змінюється з віком.

В наступному прикладі ми створюємо функцію взаємодії, яка включає в себе параметри вік та дохід. 

\begin{lstlisting}
data = {
    'age': [25, 35, 45, 55],
    'income': [30000, 50000, 70000, 90000]
}
df = pd.DataFrame(data)

df['age_income_interaction'] = df['age'] * df['income']

df
\end{lstlisting}

Результат:

\begin{console}
	age	income	age_income_interaction
0	25	30000	750000
1	35	50000	1750000
2	45	70000	3150000
3	55	90000	4950000
\end{console}



\subsection{Поліноміальні ознаки.}

Створюються шляхом генерації нових ознак, які є поліноміальними комбінаціями існуючих ознак. 
Це може бути особливо корисно при моделюванні нелінійних взаємозв'язків.

Приклад:

Припустимо, у вас є одна ознака x. Ви можете створити поліноміальні ознаки, такі як $x^2$, $x^3$ 
тощо, які можуть допомогти у виявленні нелінійних взаємозв'язків у даних.

\begin{lstlisting}
from sklearn.preprocessing import PolynomialFeatures

data = {'x': [2, 3, 4]}
df = pd.DataFrame(data)

poly = PolynomialFeatures(degree=3, include_bias=False)
poly_features = poly.fit_transform(df)

df_poly = pd.DataFrame(poly_features, columns=poly.get_feature_names_out(['x']))

df_poly
\end{lstlisting}

Результат:

\begin{console}
	x	x^2	x^3
0	2.0	4.0	8.0
1	3.0	9.0	27.0
2	4.0	16.0	64.0
\end{console}

\ding{43} Після завершення обробки необхідно переконатися, що дані не містять пропущених значень, 
дублікатів, некоректних типів даних та аномальних викидів, а всі ознаки мають очікуваний формат 
і масштаб. Для цього використовують статистичний аналіз, перевірку структури даних, візуалізацію 
та засоби бібліотеки pandas, такі як \verb|info()|, \verb|describe()|, \verb|isnull()| і 
\verb|duplicated()|. 

\ding{43} Якісно підготовлені дані дозволяють моделям машинного навчання точніше виявляти закономірності, 
швидше навчатися, бути менш чутливими до шуму та краще узагальнювати знання на нових даних, що 
безпосередньо підвищує точність, стійкість і надійність прогнозування. Особливо важливо уникати витоку інформації під час підготовки даних. Параметри імпутації, масштабування та інших перетворень потрібно визначати на навчальній вибірці, а потім застосовувати до тестових даних. Це забезпечує чесну оцінку реальної здатності моделі працювати з новими спостереженнями.

\medskip
\ding{46} \textbf{Питання для самоконтролю.}

\begin{enumerate}
\item Що таке обробка даних (Data Processing) та яке її значення в процесі інтелектуального аналізу даних?
\item Які основні етапи включає процес обробки даних?
\item Чому якість обробки даних впливає на результати машинного навчання?
\item Які основні проблеми можуть бути присутні у «сирих» даних?
\item Що таке бібліотека pandas та для чого вона використовується?
\item Які основні структури даних підтримує бібліотека pandas?
\item Що таке DataFrame?
\item Чим DataFrame відрізняється від Series?
\item Як створити DataFrame у Python?
\item Як зчитати дані у DataFrame з CSV-файлу?
\item Як переглянути перші та останні записи DataFrame?
\item Як отримати інформацію про структуру DataFrame?
\item Як отримати статистичний опис числових даних за допомогою pandas?
\item Як вибрати окремий стовпець або кілька стовпців DataFrame?
\item Яка різниця між методами `loc[]` та `iloc[]`?
\item Як відфільтрувати рядки DataFrame за певною умовою?
\item Як додати новий стовпець до DataFrame?
\item Як видалити стовпець або рядок із DataFrame?
\item Як перейменувати стовпці DataFrame?
\item Як відсортувати дані за одним або кількома стовпцями?
\item Що таке очищення даних (Data Cleaning)?
\item Які основні причини виникнення неякісних даних?
\item Що таке пропущені значення (Missing Values)?
\item Які причини появи пропущених значень у наборах даних?
\item Як виявити пропущені значення у DataFrame?
\item Які методи використовуються для видалення пропущених значень?
\item Які методи використовуються для заповнення пропущених значень?
\item У яких випадках доцільно видаляти рядки з пропущеними значеннями?
\item Що таке імпутація даних (Data Imputation)?
\item Які статистичні методи використовуються для заповнення пропущених значень?
\item Що таке дублікати даних?
\item Як знайти дублікати у DataFrame?
\item Як видалити дублікати за допомогою pandas?
\item Що таке викиди (Outliers)?
\item Якими методами можна виявляти викиди?
\item Чому викиди можуть негативно впливати на результати аналізу даних?
\item Що таке масштабування даних (Data Scaling)?
\item Для чого виконується масштабування ознак перед навчанням моделей?
\item У чому полягає різниця між масштабуванням та нормалізацією даних?
\item Що таке мінімально-максимальна нормалізація?
\item Як обчислюється мінімально-максимальна нормалізація?
\item Що таке стандартизація (Standardization)?
\item Як обчислюється стандартизація за Z-оцінкою (Z-score)?
\item Для яких алгоритмів машинного навчання масштабування є особливо важливим?
\item Які переваги має нормалізація даних?
\item Які недоліки має мінімально-максимальна нормалізація?
\item Які переваги має стандартизація порівняно з мінімально-максимальною нормалізацією?
\item Які інструменти бібліотеки scikit-learn використовуються для масштабування та нормалізації даних?
\item Як перевірити, що дані були успішно очищені та підготовлені до аналізу?
\item Яким чином правильна обробка даних впливає на точність, стійкість і узагальнювальну здатність моделей машинного навчання?
\end{enumerate}



